\section{Decomposition and Solver}\label{sec:decomposition}

\subsection{Variable Roles and Two-Level Structure}

The decision variables of \textbf{(P)} naturally split into a \emph{global/coupling block} $(\mu, \chi)$ and a \emph{link-local block} $(\beta_A, \{\beta_i\})$. Fixing the global block, the inner problem
\begin{equation}\label{eq:innerI}
\begin{aligned}
\textbf{(I)}:\quad V(\mu,\chi) = \max_{\beta_A, \{\beta_i\}}\ & \sum_{i=1}^N \rs^i \sigma_i\\[-1pt]
&\text{s.t.\ per-pair constraints of \eqref{eq:problemP}},
\end{aligned}
\end{equation}
returns a value function; the outer problem
\begin{equation}\label{eq:outerO}
\textbf{(O)}: \max_{\mu\in(0,1),\chi\in[0,1]}\ V(\mu,\chi)
\end{equation}
is then a two-dimensional search amenable to a coarse grid plus local refinement. The outer/inner \emph{split} is a lossless reformulation of \textbf{(P)}, and it is worth stating why in one line, since the two halves of the solver have very different status. Writing the feasible set as $\mathcal{F}=\{(\mu,\chi,\boldsymbol{\beta})\}$ and the objective as $\Rf(\mu,\chi,\boldsymbol{\beta})$, partial maximization gives $\max_{\mathcal{F}}\Rf = \max_{(\mu,\chi)}\big[\max_{\boldsymbol{\beta}\in\mathcal{F}(\mu,\chi)}\Rf\big]$ whenever the inner maximum exists, which it does because $\boldsymbol{\beta}$ ranges over a compact box and $\Rf$ is continuous on it. No relaxation, no bound, and no assumption beyond compactness is involved: the identity is exact. What the identity does \emph{not} give is that our inner routine attains that inner maximum. The per-user 1-D maximizations are exhaustive on a grid followed by refinement, so each attains its own grid optimum exactly; the step without a guarantee is the mean-field coupling between users, for which we have no contraction proof and therefore report empirical evidence only: agreement with a brute-force joint search at $N{=}2$, convergence within six iterations, and insensitivity to the initialization. We state the boundary between the proved and the measured explicitly rather than letting the word \emph{exact} cover both.

\subsection{Mean-Field Block-Coordinate Inner Solver}

Inner couplings inside \textbf{(I)} are of two kinds: through the shared $\beta_A$ (in every $\Pc^A, \Pe^A$), and through the exclusion field $\phi_i = \prod_{j\neq i}(1-\Pc^j)$ in \eqref{eq:excl}. We treat $\boldsymbol{\phi}$ as a field and iterate the procedure listed in Algorithm~\ref{alg:solver} and depicted in Fig.~\ref{fig:decomp}:

\begin{algorithm}[!htbp]
\caption{Decomposed solver for \textbf{(P)}}\label{alg:solver}
\begin{algorithmic}[1]
\STATE \textbf{Input:} cluster state $\bs$, grids $\mathcal{M}_\mu, \mathcal{M}_\chi$, tolerances $\epsilon_\phi, \epsilon_R$
\STATE \textbf{Output:} $\bp^\star, V^\star$
\STATE $V^\star \gets -\infty$
\FORALL{$(\mu,\chi)\in\mathcal{M}_\mu\times\mathcal{M}_\chi$}
  \STATE Init $\boldsymbol{\phi}^{(0)}\gets \boldsymbol{1}$, $\beta_A^{(0)}\gets\beta_{\mathrm{lo}}$
  \FOR{$k = 0,1,\dots$ until $\|\boldsymbol{\phi}^{(k+1)}-\boldsymbol{\phi}^{(k)}\| < \epsilon_\phi$}
    \FOR{$i = 1,\dots,N$}
      \STATE $\beta_i^\star \gets \arg\max_{\beta} \rs^i(\beta;\mu,\chi,\beta_A^{(k)},\phi_i^{(k)})\sigma_i(\cdot)$ subject to per-pair constraints \COMMENT{1-D search}
    \ENDFOR
    \STATE $\beta_A^{(k+1)} \gets \arg\max_{\beta_A} \sum_i \rs^i\sigma_i$ holding $\{\beta_i^\star\}$ \COMMENT{1-D search, block coordinate}
    \STATE $\phi_i^{(k+1)} \gets \prod_{j\neq i}(1 - \Pc^j(\beta_j^\star))$ \COMMENT{field refresh}
  \ENDFOR
  \STATE $V \gets \sum_i \rs^i \sigma_i$ at current $(\mu,\chi,\beta_A,\{\beta_i\})$
  \IF{$V > V^\star$}
    \STATE $V^\star \gets V,\quad \bp^\star \gets (\mu,\beta_A,\{\beta_i\},\chi)$
  \ENDIF
\ENDFOR
\end{algorithmic}
\end{algorithm}

The flow is depicted in Fig.~\ref{fig:decomp}.

\begin{figure}[!htbp]
\centering
\resizebox{\columnwidth}{!}{% Decomposed solver: nested-loop flowchart.
% Outer (mu,chi) grid sweep ENCLOSES the inner mean-field block-coordinate loop:
%   [per-user 1-D max over beta_i] -> [1-D max over beta_A] -> [field refresh]
%   -> convergence test (dashed feedback for "not converged").
% Emit V(mu,chi) at the inner fixed point; track global argmax p*.
\begin{tikzpicture}[
  font=\footnotesize,
  >={Latex[length=1.7mm,width=1.5mm]},
  node distance=6mm and 9mm,
  % --- shared muted palette (cohesive with the paper's other flow figures) ---
  base/.style   ={draw=black!55,rounded corners=2pt,align=center,
                  thick,text=black,inner sep=3pt},
  outerb/.style ={base,fill=blue!8,minimum width=52mm,minimum height=9mm,
                  font=\footnotesize\bfseries},
  istep/.style  ={base,fill=teal!14,minimum width=46mm,minimum height=9.5mm},
  fieldb/.style ={base,fill=black!6,minimum width=46mm,minimum height=9.5mm},
  test/.style   ={draw=black!55,diamond,aspect=2.3,fill=blue!8,text=black,
                  thick,inner sep=1pt,align=center,minimum width=22mm,
                  minimum height=11mm},
  value/.style  ={base,fill=teal!22,minimum width=46mm,minimum height=9mm},
  res/.style    ={base,fill=orange!18,draw=orange!75!black,minimum width=52mm,
                  minimum height=9.5mm,font=\footnotesize\bfseries},
  flow/.style   ={->,thick,draw=black!75},
  back/.style   ={->,thick,draw=black!55,densely dashed},
  edgelbl/.style={font=\scriptsize,text=black!60,inner sep=2pt},
]

% ============================ pipeline of nodes ============================
\node[outerb] (sweep) {Outer sweep: for each $(\mu,\chi)\in\mathcal{M}_\mu\!\times\!\mathcal{M}_\chi$};

\node[istep,below=15mm of sweep] (bi)
  {Per-user 1-D max: $\beta_i^\star=\displaystyle\arg\max_{\beta}\,\rs^i\sigma_i$\\[1pt]
   {\scriptsize $\forall i$, holding $\phi_i,\ \beta_A$ fixed}};
\node[istep,below=of bi] (ba)
  {Block-coordinate 1-D max:\\[1pt]
   $\beta_A^\star=\displaystyle\arg\max_{\beta_A}\,{\textstyle\sum_i}\rs^i\sigma_i$};
\node[fieldb,below=of ba] (field)
  {Field refresh:\\[1pt]
   $\phi_i\leftarrow\prod_{j\neq i}\!\bigl(1-\Pc^j(\beta_j^\star)\bigr)$};
\node[test,below=8mm of field] (conv) {$\phi$ \\ converged?};

\node[value,below=9mm of conv] (val)
  {Inner optimum: $V(\mu,\chi)\leftarrow{\textstyle\sum_i}\rs^i\sigma_i$};
\node[res,below=8mm of val] (res)
  {Outer argmax: $\bp^\star=\arg\max_{(\mu,\chi)} V,\quad V^\star$};

% ====================== inner-loop grouping (backgrounds) ==================
\begin{scope}[on background layer]
  \node[draw=teal!55,rounded corners=4pt,fill=teal!5,
        inner sep=4.5mm,fit=(bi)(ba)(field)(conv)] (innerbox) {};
\end{scope}
\node[anchor=south west,font=\scriptsize\itshape,text=teal!55!black,
      inner sep=1pt]
  at ($(innerbox.north west)+(2mm,0.6mm)$) {Inner: mean-field block coordinate};

% =============================== forward flow ==============================
\draw[flow] (sweep) -- node[edgelbl,right,pos=0.28]{init $\phi^{(0)}\!=\!\mathbf{1},\ \beta_A^{(0)}$} (bi);
\draw[flow] (bi)    -- (ba);
\draw[flow] (ba)    -- (field);
\draw[flow] (field) -- (conv);
\draw[flow] (conv)  -- node[edgelbl,right]{yes} (val);
\draw[flow] (val)   -- (res);

% =================== inner feedback: "not converged" =======================
% Both turn-points share x = innerbox.east + 7mm so the vertical run is exact.
\coordinate (fbx)  at ($(innerbox.east)+(7mm,0)$);
\coordinate (fbhi) at (fbx |- bi.east);
\coordinate (fblo) at (fbx |- conv.east);
\draw[back] (conv.east) -- (fblo) -- (fbhi) -- (bi.east);
\node[edgelbl,anchor=west] at ($(fblo)!0.5!(fbhi)$) {no: $k\!\leftarrow\!k\!+\!1$};

% =================== outer feedback: "next grid node" ======================
% Left side; common x = sweep.west - 10mm.
\coordinate (oux)  at ($(sweep.west)+(-10mm,0)$);
\coordinate (outop) at (oux |- sweep.west);
\coordinate (oubot) at (oux |- res.west);
\draw[back] (res.west) -- (oubot) -- (outop) -- (sweep.west);
\node[edgelbl,anchor=south,rotate=90] at ($(oubot)!0.5!(outop)+(-1.5mm,0)$) {next $(\mu,\chi)$};

\end{tikzpicture}
}
\caption{Decomposed solver. The outer grid sweeps $(\mu,\chi)$; the inner mean-field block-coordinate loop alternates per-user 1-D maximizations over $\beta_i$ (held the field $\phi_i$ and the shared $\beta_A$ fixed), then refreshes $\beta_A$ by 1-D maximization, then refreshes the field $\phi_i$ from the new $\{\Pc^j\}$. The fixed point is the inner optimum $V(\mu,\chi)$; the outer argmax is $\bp^\star$.}
\label{fig:decomp}
\end{figure}

\subsection{Correctness Checks}

\subsubsection{Symmetric-cluster reduction}
For a symmetric cluster ($s_i = s_*\,\forall i$, $\beta_i = \beta_B\,\forall i$), the field collapses to $\phi = (1-\Pc^B)^{N-1}$ and the excluded sift \eqref{eq:excl} reduces to $\Pc^A \Pc^B[1 - (1-\Pc^B)^{N-1}]$, recovering \cite[Eq.~24]{vu2023photonics}.

\subsubsection{Brute-force agreement}
On a two-user cluster the decomposition is verified against a brute-force joint grid over $(\beta_A, \beta_1, \beta_2)$ at a fixed $(\mu,\chi)$ node: the decomposed and brute-force optima agree, with a ratio of $1.000$ to four decimals (Sec.~\ref{sec:results:oracle}). The closed-form exclusion field \eqref{eq:excl} is separately checked against a direct subset-enumeration inclusion--exclusion sum for $N\in\{1,2,4,6\}$ (agreement to $10^{-12}$).

\subsection{Convergence and Complexity}

The solver uses $|\mathcal{M}_\mu| = 12$ outer nodes uniformly spaced on $[0.4, 0.95]$ and $|\mathcal{M}_\chi| = 5$ on $[0,1]$ in the headline runs. The inner field iteration converges in $\le 6$ steps from any cold start at the working tolerances $\epsilon_\phi = 10^{-3}$ and $\epsilon_R = 10^{-4}$; the iterate reached from any single initialisation converges, but the fixed point is \emph{not} unique: restarting the inner iteration from one common-$\beta$ value at a time, all $12$ of $12$ tested cluster states reach different fixed points depending on the start, with $\beta_A^\star$ differing by up to $1.31$. Multi-start is therefore not a convenience but a requirement, and the earlier draft of this paper overstated the situation by reporting uniqueness; what is unique is the best-of-four solution the solver returns, not the fixed point itself. A formal characterisation of when the map is contractive is left for future work. One full outer sweep (all $12\times5 = 60$ outer nodes, each with up to 6 inner iterations and Gauss-Hermite quadrature at $n{=}20$) takes approximately 4~s on a single CPU core, so a full pass sampled at $\Delta t = 50$~s ($\sim$50 timesteps $\times$ 4~s) solves in roughly 200~s. Multi-start is enabled for the inner 1-D maximizations, because the feasibility boundary creates several disconnected operating windows in $\beta$ for certain $(\mu, \chi)$ choices; we restart the inner iteration from four fixed common-$\beta$ values $\{0.6, 1.2, 1.8, 2.4\}$ and retain the best feasible solution, ranked by feasible-pair count and then by total rate. The starts are deterministic, so the solver is reproducible without a seed.

\subsection{Why This Decomposition Works}

Three properties make the outer / inner split practical. First, the outer block is genuinely two-dimensional: $\mu$ and $\chi$ are the only variables that couple all $N$ pairs through threats (URA via $\mu$) and through the exclusion field (via $\chi$). Second, given the outer block, the inner subproblem $(\mathrm{I})$ admits a mean-field decoupling in the user index $i$: with $\boldsymbol{\phi}$ frozen, each $\beta_i^\star$ is determined by a 1-D maximization that does not see the other users, and the field is refreshed at the next outer iteration. Third, $\beta_A$ couples back through Alice's link only, so a single 1-D block-coordinate update suffices. The whole procedure is therefore exponentially cheaper than the naive joint $(N{+}1)$-dimensional grid search, while preserving exactness in the symmetric special case (Sec.~\ref{sec:results:oracle}).
